\section{结论}

\subsection{主要发现}

本研究以药物母体为评价单位，比较了描述符扩展和异质模型融合对条件溶解度预测的影响。基础 ExtraTrees、丰富特征 ExtraTrees 和双分支融合在开发集的母体等权 MAE 分别为 0.7469、0.6831 和 0.6668。主要改善来自树分支整体设计的变化，融合在当前比较中提供了较小的增量，且配对区间反映出该增量仍有波动。误差分析将低溶解度端确定为优先核验的场景。

\subsection{方法与应用交付}

交付包括从结构描述符计算、训练折内选择到母体折外评价的建模流程，以及部署双分支模型的交互应用。系统可在指定条件下给出预测、分支贡献、温度曲线和开发集覆盖信息。后续可将固定模型用于独立来源数据，先检查低溶解度端和新化学骨架上的偏差，再评估校准方法能否改善筛选表现。
